\section{A learned pose initialiser, and where it breaks}
\label{sec:pointnet-init}

\S\ref{sec:init-stage} introduced a \term{PointNet}-style~\cite{pointnet}
initialiser: a point-cloud encoder that reads the raw target points,
extracts a permutation-invariant feature per point, pools them into one
global descriptor, and regresses an initial pose and shape directly from
it, replacing a fixed canonical starting pose with a learned, input-aware
one. The motivation follows from \F{F4}: for the distal chain, initialisation determines
which target surface region becomes each vertex's first attractor, so a
better starting point should, in principle, reduce how often the optimiser
locks onto the wrong one.

\begin{figure}[!htb]
\centering
\begin{tikzpicture}[every node/.style={font=\scriptsize},
 box/.style={draw,rounded corners=2pt,minimum height=8mm,align=center,fill=white},
 arr/.style={-{Stealth[length=1.6mm]}}]
\node[box] (a) {point cloud\\$N\times3$};
\node[box,right=6mm of a] (b) {shared MLP\\per point};
\node[box,right=6mm of b] (c) {symmetric\\max-pool};
\node[box,right=6mm of c] (d) {global\\descriptor};
\node[box,right=6mm of d] (e) {regression head\\$\hat{\bm\theta}$};
\node[box,right=6mm of e] (f) {SMILify\\optimiser};
\foreach \x/\y in {a/b,b/c,c/d,d/e,e/f}{\draw[arr] (\x)--(\y);}
\end{tikzpicture}
\caption{Learned initialiser. Max-pooling makes the descriptor invariant to
point ordering, which a raw scan does not have.}
\label{fig:pointnet}
\end{figure}


\begin{investigation}{I11}{Does the learned initialiser improve registration,
and does a synthetic improvement transfer to real specimens?}
\invQuestion{Under controlled, exact synthetic ground truth, does the
learned initialiser outperform the existing default; and does that
advantage hold on real scans?}
\invMethod{Synthetic benchmark: learned vs.\ default initialisation, exact
ground truth known. Real-data benchmark: 50 real specimens, learned vs.\
zero (default) initialisation at three levels of synthetic pose
perturbation used during training ($\approx\!23^\circ$, then retrained at
$\approx\!13.8^\circ$ and $\approx\!8.85^\circ$).}
\invResult{On synthetic data the learned initialiser wins clearly: leg
accuracy improves from $0.857$ to $0.901$ and Chamfer distance from
$0.000264$ to $0.000226$. On the 50-specimen real benchmark at the
$23^\circ$ training regime, this reverses: learned initialisation
($0.000697$ mean Chamfer) is \emph{worse} than the zero-initialised default
($0.000637$), winning only 20 of 50 specimens. Retraining at smaller
perturbation magnitudes moves the comparison toward parity rather than
robust superiority.}
\invVerdict{REJECT}{The observed failure is consistent with a
synthetic-to-real distribution mismatch: performance did not transfer from
the tested synthetic training regime to real scans, and the
perturbation-magnitude sweep supports that interpretation. The experiments
do not establish whether the architecture itself would generalise under a
more representative training distribution, only that it did not under the
one tested. Not deployed in production as evaluated. The general lesson is
retained as a standing caution: a better initialiser can, in principle,
improve optimisation without addressing correspondence at all, and a
synthetic evaluation cannot certify real-world transfer by itself.}
\end{investigation}

A controlled follow-up sweep, holding total pose-error magnitude fixed and
varying only where along the kinematic chain it was placed, is what produced
the proximal/distal result reported as \F{F4} in \S\ref{sec:articulation}:
the same error placed proximally degrades leg accuracy far more (to $0.768$)
than the identical error placed distally ($0.939$), because a proximal
mistake propagates through every joint downstream of it.
